% Part: second-order-logic
% Chapter: syntax-and-semantics
% Section: introduction

\documentclass[../../../include/open-logic-section]{subfiles}

\begin{document}

\olfileid{sol}{syn}{int}

\olsection{పరిచయం}

మొదటిస్థాయి తర్కంలో, ఇచ్చిన భాషలోని తార్కికేతర సంకేతాలను, అంటే దాని
\tetoken{స్థిరసంకేతాలు}{!!{constant}s}, \tetoken{ప్రమేయ సంకేతాలు}{!!{function}s}, \tetoken{విధేయ సంకేతాలను}{!!{predicate}s}, తార్కిక
సంకేతాలతో కలిపి మొదటిస్థాయి \tetoken{నిర్మాణాల}{!!{structure}s} గురించి విషయాలను వ్యక్తం
చేస్తాం. ఇందుకు సంతృప్తి భావనను ఉపయోగిస్తాం; అది
\tetoken{ఒక నిర్మాణం}{!!a{structure}}~$\Struct{M}$, ఒక చర నిర్దేశం~$s$, \tetoken{ఒక సూత్రం}{!!a{formula}}~$!A$ల
మధ్య సంబంధాన్ని ఏర్పరుస్తుంది: $!A$లోని \tetoken{స్థిరసంకేతాలు}{!!{constant}s},
\tetoken{ప్రమేయ సంకేతాలు}{!!{function}s}, \tetoken{విధేయ సంకేతాలను}{!!{predicate}s} $\Struct{M}$ చెప్పినట్లుగా,
దాని స్వేచ్ఛా చరాలను~$s$ చెప్పినట్లుగా అర్థనిర్దేశం చేసినప్పుడు వ్యక్తమైనది
సత్యమైతే, అప్పుడు మాత్రమే $\Sat{M}{!A}[s]$. $\eq$ అనే
\tetoken{తాదాత్మ్య విధేయం}{!!{identity}} యొక్క అర్థనిర్దేశం, అలాగే $\lforall$, $\lexists$ల
అర్థనిర్దేశం $\Sat{M}{!A}[s]$ నిర్వచనంలోనే అంతర్నిర్మితంగా ఉంటాయి.
మొదటిదాన్ని నిర్మాణపు \tetoken{వ్యక్తి క్షేత్రం}{!!{domain}}~$\Domain{M}$పై తాదాత్మ్య
సంబంధంగానే ఎల్లప్పుడూ అర్థనిర్దేశం చేస్తాం; పరిమాణకాలు ఎల్లప్పుడూ
మొత్తం \tetoken{వ్యక్తి క్షేత్రం}{!!{domain}}పై విస్తరించినట్లుగా అర్థనిర్దేశం చేస్తాం. అయితే
కీలకంగా, పరిమాణీకరణను \tetoken{వ్యక్తి క్షేత్రంలోని}{!!{domain}} మూలకాలపైనే అనుమతిస్తాం;
అందువల్ల పరిమాణకం తరువాత వస్తు \tetoken{చరాలను}{!!{variable}s} మాత్రమే అనుమతిస్తాం.

ద్వితీయ-స్థాయి తర్కంలో, స్వేచ్ఛా, బద్ధ ప్రమేయ చరాలు, విధేయ చరాలు,
వాటిపై పరిమాణీకరణను చేర్చేలా భాషనూ సంతృప్తి నిర్వచనాన్నీ
విస్తరిస్తాం. వస్తు చరాలకు \tetoken{స్థిరసంకేతాలతో}{!!{constant}s} ఎలాంటి సంబంధమో, ఈ చరాలకు
\tetoken{ప్రమేయ సంకేతాలు}{!!{function}s}, \tetoken{విధేయ సంకేతాలతో}{!!{predicate}s} అలాంటి సంబంధమే ఉంటుంది.
ద్వితీయ-స్థాయి తర్కపు పదాలు, \tetoken{సూత్రాల}{!!{formula}s} నిర్మాణంలో అవి అదే పాత్ర
పోషిస్తాయి; వాటిపై పరిమాణీకరణను కూడా అదే విధంగా నిర్వహిస్తాం.
\emph{ప్రామాణిక} అర్థవిచారంలో, ద్వితీయ-స్థాయి పరిమాణకాలు సరైన రకానికి
చెందిన సాధ్యమైన వస్తువులన్నింటిపై విస్తరిస్తాయి (ప్రమేయ చరాల కోసం
$\Domain{M}$ నుంచి~$\Domain{M}$కు $n$-స్థాన ప్రమేయాలు, విధేయ చరాల
కోసం $n$-స్థాన సంబంధాలు). ఉదాహరణకు,
$\lforall[\Obj{v_0}][(\Obj{P^1_0}(\Obj{v_0}) \lor \lnot
  \Obj{P^1_0}(\Obj{v_0}))]$ అనేది మొదటిస్థాయి, ద్వితీయ-స్థాయి
తర్కాలు రెండింటిలోనూ ఒక సూత్రం; రెండవదానిలో మనం
$\lforall[\Obj{V^1_0}][\lforall[\Obj{v_0}][(\Obj{V^1_0}(\Obj{v_0}) \lor
    \lnot \Obj{V^1_0}(\Obj{v_0}))]]$, అలాగే
$\lexists[\Obj{V^1_0}][\lforall[\Obj{v_0}][(\Obj{V^1_0}(\Obj{v_0}) \lor
    \lnot \Obj{V^1_0}(\Obj{v_0}))]]$లను కూడా పరిగణించవచ్చు. వీటిలో
స్వేచ్ఛా చరాలు లేవు కాబట్టి, అవి ద్వితీయ-స్థాయి తర్కపు
\tetoken{వాక్యాలు}{!!{sentence}s}. ఇక్కడ $\Obj{V^1_0}$ ఒక $1$-స్థాన ద్వితీయ-స్థాయి విధేయ
చరం. $\Obj{V^1_0}$కు అనుమతించే అర్థనిర్దేశాలు, $\Obj{P^1_0}$ వంటి
$1$-స్థాన \tetoken{విధేయ సంకేతానికి}{!!{predicate}} మనం నిర్దేశించగల వాటితో సమానమైనవి, అంటే
అవి~$\Domain{M}$ యొక్క ఉపసమితులు. వాటిపై పరిమాణీకరించడం అంటే,
$\Domain{M}$ యొక్క ఒక ఉపసమితిని $\Obj{V^1_0}$ విలువగా నిర్దేశించే
అన్ని విధానాలకూ, లేదా కనీసం ఒక విధానానికైనా,
$\lforall[\Obj{v_0}][(\Obj{V^1_0}(\Obj v_0) \lor \lnot
  \Obj{V^1_0}(v_0))]$ సత్యమని చెప్పడమే. ప్రతి సమితి ఇచ్చిన వస్తువును
కలిగి ఉంటుంది లేదా కలిగి ఉండదు కాబట్టి, రెండూ ఏ
\tetoken{నిర్మాణంలోనైనా}{!!{structure}} సత్యమే.
\end{document}
